% Подключение расчетного модуля Lua
\directlua{dofile('processing.lua')}

\begin{figure}[hb]
    \centering
    % Здесь должен быть график, построенный на основе данных.
    % В данном случае используется заглушка.
    \includegraphics[width=\textwidth]{figs/stress-strain-plot.png}
    \caption{Диаграммы условных напряжений для исследованных образцов}
    \label{fig:stress_strain_src}
\end{figure}

\section*{ \boxed{\text{ Задание 1. }} }
\begin{quote}
    \textit{Вычисление площадей поперечного сечения до ($F_0$) и после ($F_k$) испытания.}
\end{quote}

Площадь поперечного сечения для образцов с прямоугольным профилем (кожа) вычисляется по формуле:
\begin{equation*}
    F = w \cdot h,
\end{equation*}
где $w$ --- ширина, $h$ --- толщина образца.

Для образцов с круглым профилем (веревка, проволока) формула имеет вид:
\begin{equation*}
    F = \frac{\pi d^2}{4},
\end{equation*}
где $d$ --- диаметр образца.

Результаты расчетов сведены в \cref{tab:areas}.

\begin{table}[H]
    \centering
    \caption{Расчетные значения площадей поперечного сечения}
    \label{tab:areas}
    \begin{tblr}{
        colspec = { Q[l,m] *{2}{X[c,m]} },
    }
        \toprule
        Образец & Начальная площадь $F_0$, м$^2$ & Конечная площадь $F_k$, м$^2$ \\
        \midrule
        Кожа     & \directlua{format_sci(results.leather.F0)} & \directlua{format_sci(results.leather.Fk)} \\
        Веревка  & \directlua{format_sci(results.rope.F0)}    & \directlua{format_sci(results.rope.Fk)}    \\
        Проволока& \directlua{format_sci(results.wire.F0)}    & \directlua{format_sci(results.wire.Fk)}    \\
        \bottomrule
    \end{tblr}
\end{table}

\section*{ \boxed{\text{ Задание 2. }} }
\begin{quote}
    \textit{Расчет механических характеристик прочности ($\sigma_\text{в}, \sigma_\text{к}$).}
\end{quote}

Для расчета характеристик прочности используются значения нагрузок, снятые с диаграммы растяжения. Максимальная нагрузка $P_\text{max}$ соответствует временному сопротивлению материала, а нагрузка в момент разрыва $P_k$ --- истинному сопротивлению разрыву.

\begin{table}[H]
    \centering
    \caption{Значения нагрузок, определенные по диаграмме растяжения}
    \label{tab:loads}
    \begin{tblr}{
        colspec = { Q[l,m] *{2}{X[c,m]} },
    }
        \toprule
        Образец & $P_\text{max}$, Н & $P_k$, Н \\
        \midrule
        Кожа     & \directlua{tex.print(P_values.leather.max)} & \directlua{tex.print(P_values.leather.k)} \\
        Веревка  & \directlua{tex.print(P_values.rope.max)}    & \directlua{tex.print(P_values.rope.k)} \\
        Проволока& \directlua{tex.print(P_values.wire.max)}    & \directlua{tex.print(P_values.wire.k)} \\
        \bottomrule
    \end{tblr}
\end{table}

\textit{Временное сопротивление (предел прочности)} $\sigma_\text{в}$ определяется как отношение максимальной нагрузки к начальной площади поперечного сечения:
\begin{equation}
    \sigma_\text{в} = \frac{P_\text{max}}{F_0}.
    \label{eq:sigma_v}
\end{equation}

\textit{Истинное сопротивление разрыву} $\sigma_\text{к}$ определяется как отношение разрывной нагрузки к площади сечения шейки в момент разрыва:
\begin{equation}
    \sigma_\text{к} = \frac{P_k}{F_k}.
    \label{eq:sigma_k}
\end{equation}

Пример расчета для образца <<Проволока>>:
\begin{equation*}
    \sigma_\text{в} = \frac{\val[0]{P_values.wire.max} \text{ Н}}{\val[exp]{results.wire.F0} \text{ м}^2} \approx \val[exp]{results.wire.sigma_v} \text{ Па}.
\end{equation*}
\begin{equation*}
    \sigma_\text{к} = \frac{\val[0]{P_values.wire.k} \text{ Н}}{\val[exp]{results.wire.Fk} \text{ м}^2} \approx \val[exp]{results.wire.sigma_k} \text{ Па}.
\end{equation*}

Результаты расчетов для всех образцов приведены в \cref{tab:strength}.

\begin{table}[H]
    \centering
    \caption{Характеристики прочности материалов}
    \label{tab:strength}
    \begin{tblr}{
        colspec = { Q[l,m] *{2}{X[c,m]} },
    }
        \toprule
        Образец & Временное сопротивление $\sigma_\text{в}$, МПа & Истинное сопротивление разрыву $\sigma_\text{к}$, МПа \\
        \midrule
        Кожа     & \directlua{cell(results.leather.sigma_v / 1e6)} & \directlua{cell(results.leather.sigma_k / 1e6)} \\
        Веревка  & \directlua{cell(results.rope.sigma_v / 1e6)}    & \directlua{cell(results.rope.sigma_k / 1e6)}    \\
        Проволока& \directlua{cell(results.wire.sigma_v / 1e6)}    & \directlua{cell(results.wire.sigma_k / 1e6)}    \\
        \bottomrule
    \end{tblr}
\end{table}


\section*{ \boxed{\text{ Задание 3. }} }
\begin{quote}
    \textit{Расчет характеристик пластичности ($\delta, \psi$).}
\end{quote}

\textit{Относительное удлинение после разрыва} $\delta$ вычисляется по формуле:
\begin{equation}
    \delta = \frac{l_k - l_0}{l_0} \cdot 100\%,
    \label{eq:delta}
\end{equation}
где $l_0$ и $l_k$ --- начальная и конечная расчетная длина образца.

\textit{Относительное сужение после разрыва} $\psi$ вычисляется по формуле:
\begin{equation}
    \psi = \frac{F_0 - F_k}{F_0} \cdot 100\%.
    \label{eq:psi}
\end{equation}

Результаты расчетов сведены в \cref{tab:plasticity}.

\begin{table}[H]
    \centering
    \caption{Характеристики пластичности материалов}
    \label{tab:plasticity}
    \begin{tblr}{
        colspec = { Q[l,m] *{2}{X[c,m]} },
    }
        \toprule
        Образец & Относительное удлинение $\delta$, \% & Относительное сужение $\psi$, \% \\
        \midrule
        Кожа     & \directlua{cell(results.leather.delta)} & \directlua{cell(results.leather.psi)} \\
        Веревка  & \directlua{cell(results.rope.delta)}    & \directlua{cell(results.rope.psi)}    \\
        Проволока& \directlua{cell(results.wire.delta)}    & \directlua{cell(results.wire.psi)}    \\
        \bottomrule
    \end{tblr}
\end{table}

\section*{ \boxed{\text{ Задание 4. }} }
\begin{quote}
    \textit{Построение диаграммы условных напряжений $\sigma = f(\varepsilon)$.}
\end{quote}

Диаграмма условных напряжений строится путем пересчета осей диаграммы растяжения <<нагрузка-удлинение>>. Ось ординат (напряжение $\sigma$) получается делением нагрузки $P$ на начальную площадь $F_0$. Ось абсцисс (относительная деформация $\varepsilon$) --- делением абсолютного удлинения $\Delta l$ на начальную длину $l_0$.
\begin{equation*}
    \sigma = \frac{P}{F_0}, \quad \varepsilon = \frac{\Delta l}{l_0}.
\end{equation*}
Качественный вид диаграммы $\sigma(\varepsilon)$ повторяет вид исходной диаграммы $P(\Delta l)$, представленной в условии. На \cref{fig:stress_strain} показан результат такого пересчета.
 
\begin{figure}[H]
    \centering
    % Здесь должен быть график, построенный на основе данных.
    % В данном случае используется заглушка.
    %% Creator: Matplotlib, PGF backend
%%
%% To include the figure in your LaTeX document, write
%%   \input{<filename>.pgf}
%%
%% Make sure the required packages are loaded in your preamble
%%   \usepackage{pgf}
%%
%% Also ensure that all the required font packages are loaded; for instance,
%% the lmodern package is sometimes necessary when using math font.
%%   \usepackage{lmodern}
%%
%% Figures using additional raster images can only be included by \input if
%% they are in the same directory as the main LaTeX file. For loading figures
%% from other directories you can use the `import` package
%%   \usepackage{import}
%%
%% and then include the figures with
%%   \import{<path to file>}{<filename>.pgf}
%%
%% Matplotlib used the following preamble
%%   \def\mathdefault#1{#1}
%%   \everymath=\expandafter{\the\everymath\displaystyle}
%%   \IfFileExists{scrextend.sty}{
%%     \usepackage[fontsize=14.000000pt]{scrextend}
%%   }{
%%     \renewcommand{\normalsize}{\fontsize{14.000000}{16.800000}\selectfont}
%%     \normalsize
%%   }
%%   \usepackage{fontspec}\setmainfont{Times New Roman}\usepackage[english,russian]{babel}\usepackage{amsmath}
%%   \ifdefined\pdftexversion\else  % non-pdftex case.
%%     \usepackage{fontspec}
%%     \setmainfont{times.ttf}[Path=\detokenize{/usr/local/share/fonts/truetype/times-new-roman/}]
%%     \setsansfont{DejaVuSans.ttf}[Path=\detokenize{/usr/share/matplotlib/mpl-data/fonts/ttf/}]
%%     \setmonofont{DejaVuSansMono.ttf}[Path=\detokenize{/usr/share/matplotlib/mpl-data/fonts/ttf/}]
%%   \fi
%%   \makeatletter\@ifpackageloaded{underscore}{}{\usepackage[strings]{underscore}}\makeatother
%%
\begingroup%
\makeatletter%
\begin{pgfpicture}%
\pgfpathrectangle{\pgfpointorigin}{\pgfqpoint{5.680000in}{3.680000in}}%
\pgfusepath{use as bounding box, clip}%
\begin{pgfscope}%
\pgfsetbuttcap%
\pgfsetmiterjoin%
\definecolor{currentfill}{rgb}{1.000000,1.000000,1.000000}%
\pgfsetfillcolor{currentfill}%
\pgfsetlinewidth{0.000000pt}%
\definecolor{currentstroke}{rgb}{1.000000,1.000000,1.000000}%
\pgfsetstrokecolor{currentstroke}%
\pgfsetdash{}{0pt}%
\pgfpathmoveto{\pgfqpoint{0.000000in}{0.000000in}}%
\pgfpathlineto{\pgfqpoint{5.680000in}{0.000000in}}%
\pgfpathlineto{\pgfqpoint{5.680000in}{3.680000in}}%
\pgfpathlineto{\pgfqpoint{0.000000in}{3.680000in}}%
\pgfpathlineto{\pgfqpoint{0.000000in}{0.000000in}}%
\pgfpathclose%
\pgfusepath{fill}%
\end{pgfscope}%
\begin{pgfscope}%
\pgfsetbuttcap%
\pgfsetmiterjoin%
\definecolor{currentfill}{rgb}{1.000000,1.000000,1.000000}%
\pgfsetfillcolor{currentfill}%
\pgfsetlinewidth{0.000000pt}%
\definecolor{currentstroke}{rgb}{0.000000,0.000000,0.000000}%
\pgfsetstrokecolor{currentstroke}%
\pgfsetstrokeopacity{0.000000}%
\pgfsetdash{}{0pt}%
\pgfpathmoveto{\pgfqpoint{0.624162in}{0.530740in}}%
\pgfpathlineto{\pgfqpoint{5.630000in}{0.530740in}}%
\pgfpathlineto{\pgfqpoint{5.630000in}{3.630000in}}%
\pgfpathlineto{\pgfqpoint{0.624162in}{3.630000in}}%
\pgfpathlineto{\pgfqpoint{0.624162in}{0.530740in}}%
\pgfpathclose%
\pgfusepath{fill}%
\end{pgfscope}%
\begin{pgfscope}%
\pgfpathrectangle{\pgfqpoint{0.624162in}{0.530740in}}{\pgfqpoint{5.005838in}{3.099260in}}%
\pgfusepath{clip}%
\pgfsetrectcap%
\pgfsetroundjoin%
\pgfsetlinewidth{0.501875pt}%
\definecolor{currentstroke}{rgb}{0.501961,0.501961,0.501961}%
\pgfsetstrokecolor{currentstroke}%
\pgfsetstrokeopacity{0.500000}%
\pgfsetdash{}{0pt}%
\pgfpathmoveto{\pgfqpoint{0.624162in}{0.530740in}}%
\pgfpathlineto{\pgfqpoint{0.624162in}{3.630000in}}%
\pgfusepath{stroke}%
\end{pgfscope}%
\begin{pgfscope}%
\pgfsetbuttcap%
\pgfsetroundjoin%
\definecolor{currentfill}{rgb}{0.000000,0.000000,0.000000}%
\pgfsetfillcolor{currentfill}%
\pgfsetlinewidth{0.803000pt}%
\definecolor{currentstroke}{rgb}{0.000000,0.000000,0.000000}%
\pgfsetstrokecolor{currentstroke}%
\pgfsetdash{}{0pt}%
\pgfsys@defobject{currentmarker}{\pgfqpoint{0.000000in}{-0.048611in}}{\pgfqpoint{0.000000in}{0.000000in}}{%
\pgfpathmoveto{\pgfqpoint{0.000000in}{0.000000in}}%
\pgfpathlineto{\pgfqpoint{0.000000in}{-0.048611in}}%
\pgfusepath{stroke,fill}%
}%
\begin{pgfscope}%
\pgfsys@transformshift{0.624162in}{0.530740in}%
\pgfsys@useobject{currentmarker}{}%
\end{pgfscope}%
\end{pgfscope}%
\begin{pgfscope}%
\definecolor{textcolor}{rgb}{0.000000,0.000000,0.000000}%
\pgfsetstrokecolor{textcolor}%
\pgfsetfillcolor{textcolor}%
\pgftext[x=0.624162in,y=0.433518in,,top]{\color{textcolor}{\rmfamily\fontsize{12.000000}{14.400000}\selectfont\catcode`\^=\active\def^{\ifmmode\sp\else\^{}\fi}\catcode`\%=\active\def%{\%}$\mathdefault{0}$}}%
\end{pgfscope}%
\begin{pgfscope}%
\pgfpathrectangle{\pgfqpoint{0.624162in}{0.530740in}}{\pgfqpoint{5.005838in}{3.099260in}}%
\pgfusepath{clip}%
\pgfsetrectcap%
\pgfsetroundjoin%
\pgfsetlinewidth{0.501875pt}%
\definecolor{currentstroke}{rgb}{0.501961,0.501961,0.501961}%
\pgfsetstrokecolor{currentstroke}%
\pgfsetstrokeopacity{0.500000}%
\pgfsetdash{}{0pt}%
\pgfpathmoveto{\pgfqpoint{1.540982in}{0.530740in}}%
\pgfpathlineto{\pgfqpoint{1.540982in}{3.630000in}}%
\pgfusepath{stroke}%
\end{pgfscope}%
\begin{pgfscope}%
\pgfsetbuttcap%
\pgfsetroundjoin%
\definecolor{currentfill}{rgb}{0.000000,0.000000,0.000000}%
\pgfsetfillcolor{currentfill}%
\pgfsetlinewidth{0.803000pt}%
\definecolor{currentstroke}{rgb}{0.000000,0.000000,0.000000}%
\pgfsetstrokecolor{currentstroke}%
\pgfsetdash{}{0pt}%
\pgfsys@defobject{currentmarker}{\pgfqpoint{0.000000in}{-0.048611in}}{\pgfqpoint{0.000000in}{0.000000in}}{%
\pgfpathmoveto{\pgfqpoint{0.000000in}{0.000000in}}%
\pgfpathlineto{\pgfqpoint{0.000000in}{-0.048611in}}%
\pgfusepath{stroke,fill}%
}%
\begin{pgfscope}%
\pgfsys@transformshift{1.540982in}{0.530740in}%
\pgfsys@useobject{currentmarker}{}%
\end{pgfscope}%
\end{pgfscope}%
\begin{pgfscope}%
\definecolor{textcolor}{rgb}{0.000000,0.000000,0.000000}%
\pgfsetstrokecolor{textcolor}%
\pgfsetfillcolor{textcolor}%
\pgftext[x=1.540982in,y=0.433518in,,top]{\color{textcolor}{\rmfamily\fontsize{12.000000}{14.400000}\selectfont\catcode`\^=\active\def^{\ifmmode\sp\else\^{}\fi}\catcode`\%=\active\def%{\%}$\mathdefault{10}$}}%
\end{pgfscope}%
\begin{pgfscope}%
\pgfpathrectangle{\pgfqpoint{0.624162in}{0.530740in}}{\pgfqpoint{5.005838in}{3.099260in}}%
\pgfusepath{clip}%
\pgfsetrectcap%
\pgfsetroundjoin%
\pgfsetlinewidth{0.501875pt}%
\definecolor{currentstroke}{rgb}{0.501961,0.501961,0.501961}%
\pgfsetstrokecolor{currentstroke}%
\pgfsetstrokeopacity{0.500000}%
\pgfsetdash{}{0pt}%
\pgfpathmoveto{\pgfqpoint{2.457802in}{0.530740in}}%
\pgfpathlineto{\pgfqpoint{2.457802in}{3.630000in}}%
\pgfusepath{stroke}%
\end{pgfscope}%
\begin{pgfscope}%
\pgfsetbuttcap%
\pgfsetroundjoin%
\definecolor{currentfill}{rgb}{0.000000,0.000000,0.000000}%
\pgfsetfillcolor{currentfill}%
\pgfsetlinewidth{0.803000pt}%
\definecolor{currentstroke}{rgb}{0.000000,0.000000,0.000000}%
\pgfsetstrokecolor{currentstroke}%
\pgfsetdash{}{0pt}%
\pgfsys@defobject{currentmarker}{\pgfqpoint{0.000000in}{-0.048611in}}{\pgfqpoint{0.000000in}{0.000000in}}{%
\pgfpathmoveto{\pgfqpoint{0.000000in}{0.000000in}}%
\pgfpathlineto{\pgfqpoint{0.000000in}{-0.048611in}}%
\pgfusepath{stroke,fill}%
}%
\begin{pgfscope}%
\pgfsys@transformshift{2.457802in}{0.530740in}%
\pgfsys@useobject{currentmarker}{}%
\end{pgfscope}%
\end{pgfscope}%
\begin{pgfscope}%
\definecolor{textcolor}{rgb}{0.000000,0.000000,0.000000}%
\pgfsetstrokecolor{textcolor}%
\pgfsetfillcolor{textcolor}%
\pgftext[x=2.457802in,y=0.433518in,,top]{\color{textcolor}{\rmfamily\fontsize{12.000000}{14.400000}\selectfont\catcode`\^=\active\def^{\ifmmode\sp\else\^{}\fi}\catcode`\%=\active\def%{\%}$\mathdefault{20}$}}%
\end{pgfscope}%
\begin{pgfscope}%
\pgfpathrectangle{\pgfqpoint{0.624162in}{0.530740in}}{\pgfqpoint{5.005838in}{3.099260in}}%
\pgfusepath{clip}%
\pgfsetrectcap%
\pgfsetroundjoin%
\pgfsetlinewidth{0.501875pt}%
\definecolor{currentstroke}{rgb}{0.501961,0.501961,0.501961}%
\pgfsetstrokecolor{currentstroke}%
\pgfsetstrokeopacity{0.500000}%
\pgfsetdash{}{0pt}%
\pgfpathmoveto{\pgfqpoint{3.374622in}{0.530740in}}%
\pgfpathlineto{\pgfqpoint{3.374622in}{3.630000in}}%
\pgfusepath{stroke}%
\end{pgfscope}%
\begin{pgfscope}%
\pgfsetbuttcap%
\pgfsetroundjoin%
\definecolor{currentfill}{rgb}{0.000000,0.000000,0.000000}%
\pgfsetfillcolor{currentfill}%
\pgfsetlinewidth{0.803000pt}%
\definecolor{currentstroke}{rgb}{0.000000,0.000000,0.000000}%
\pgfsetstrokecolor{currentstroke}%
\pgfsetdash{}{0pt}%
\pgfsys@defobject{currentmarker}{\pgfqpoint{0.000000in}{-0.048611in}}{\pgfqpoint{0.000000in}{0.000000in}}{%
\pgfpathmoveto{\pgfqpoint{0.000000in}{0.000000in}}%
\pgfpathlineto{\pgfqpoint{0.000000in}{-0.048611in}}%
\pgfusepath{stroke,fill}%
}%
\begin{pgfscope}%
\pgfsys@transformshift{3.374622in}{0.530740in}%
\pgfsys@useobject{currentmarker}{}%
\end{pgfscope}%
\end{pgfscope}%
\begin{pgfscope}%
\definecolor{textcolor}{rgb}{0.000000,0.000000,0.000000}%
\pgfsetstrokecolor{textcolor}%
\pgfsetfillcolor{textcolor}%
\pgftext[x=3.374622in,y=0.433518in,,top]{\color{textcolor}{\rmfamily\fontsize{12.000000}{14.400000}\selectfont\catcode`\^=\active\def^{\ifmmode\sp\else\^{}\fi}\catcode`\%=\active\def%{\%}$\mathdefault{30}$}}%
\end{pgfscope}%
\begin{pgfscope}%
\pgfpathrectangle{\pgfqpoint{0.624162in}{0.530740in}}{\pgfqpoint{5.005838in}{3.099260in}}%
\pgfusepath{clip}%
\pgfsetrectcap%
\pgfsetroundjoin%
\pgfsetlinewidth{0.501875pt}%
\definecolor{currentstroke}{rgb}{0.501961,0.501961,0.501961}%
\pgfsetstrokecolor{currentstroke}%
\pgfsetstrokeopacity{0.500000}%
\pgfsetdash{}{0pt}%
\pgfpathmoveto{\pgfqpoint{4.291443in}{0.530740in}}%
\pgfpathlineto{\pgfqpoint{4.291443in}{3.630000in}}%
\pgfusepath{stroke}%
\end{pgfscope}%
\begin{pgfscope}%
\pgfsetbuttcap%
\pgfsetroundjoin%
\definecolor{currentfill}{rgb}{0.000000,0.000000,0.000000}%
\pgfsetfillcolor{currentfill}%
\pgfsetlinewidth{0.803000pt}%
\definecolor{currentstroke}{rgb}{0.000000,0.000000,0.000000}%
\pgfsetstrokecolor{currentstroke}%
\pgfsetdash{}{0pt}%
\pgfsys@defobject{currentmarker}{\pgfqpoint{0.000000in}{-0.048611in}}{\pgfqpoint{0.000000in}{0.000000in}}{%
\pgfpathmoveto{\pgfqpoint{0.000000in}{0.000000in}}%
\pgfpathlineto{\pgfqpoint{0.000000in}{-0.048611in}}%
\pgfusepath{stroke,fill}%
}%
\begin{pgfscope}%
\pgfsys@transformshift{4.291443in}{0.530740in}%
\pgfsys@useobject{currentmarker}{}%
\end{pgfscope}%
\end{pgfscope}%
\begin{pgfscope}%
\definecolor{textcolor}{rgb}{0.000000,0.000000,0.000000}%
\pgfsetstrokecolor{textcolor}%
\pgfsetfillcolor{textcolor}%
\pgftext[x=4.291443in,y=0.433518in,,top]{\color{textcolor}{\rmfamily\fontsize{12.000000}{14.400000}\selectfont\catcode`\^=\active\def^{\ifmmode\sp\else\^{}\fi}\catcode`\%=\active\def%{\%}$\mathdefault{40}$}}%
\end{pgfscope}%
\begin{pgfscope}%
\pgfpathrectangle{\pgfqpoint{0.624162in}{0.530740in}}{\pgfqpoint{5.005838in}{3.099260in}}%
\pgfusepath{clip}%
\pgfsetrectcap%
\pgfsetroundjoin%
\pgfsetlinewidth{0.501875pt}%
\definecolor{currentstroke}{rgb}{0.501961,0.501961,0.501961}%
\pgfsetstrokecolor{currentstroke}%
\pgfsetstrokeopacity{0.500000}%
\pgfsetdash{}{0pt}%
\pgfpathmoveto{\pgfqpoint{5.208263in}{0.530740in}}%
\pgfpathlineto{\pgfqpoint{5.208263in}{3.630000in}}%
\pgfusepath{stroke}%
\end{pgfscope}%
\begin{pgfscope}%
\pgfsetbuttcap%
\pgfsetroundjoin%
\definecolor{currentfill}{rgb}{0.000000,0.000000,0.000000}%
\pgfsetfillcolor{currentfill}%
\pgfsetlinewidth{0.803000pt}%
\definecolor{currentstroke}{rgb}{0.000000,0.000000,0.000000}%
\pgfsetstrokecolor{currentstroke}%
\pgfsetdash{}{0pt}%
\pgfsys@defobject{currentmarker}{\pgfqpoint{0.000000in}{-0.048611in}}{\pgfqpoint{0.000000in}{0.000000in}}{%
\pgfpathmoveto{\pgfqpoint{0.000000in}{0.000000in}}%
\pgfpathlineto{\pgfqpoint{0.000000in}{-0.048611in}}%
\pgfusepath{stroke,fill}%
}%
\begin{pgfscope}%
\pgfsys@transformshift{5.208263in}{0.530740in}%
\pgfsys@useobject{currentmarker}{}%
\end{pgfscope}%
\end{pgfscope}%
\begin{pgfscope}%
\definecolor{textcolor}{rgb}{0.000000,0.000000,0.000000}%
\pgfsetstrokecolor{textcolor}%
\pgfsetfillcolor{textcolor}%
\pgftext[x=5.208263in,y=0.433518in,,top]{\color{textcolor}{\rmfamily\fontsize{12.000000}{14.400000}\selectfont\catcode`\^=\active\def^{\ifmmode\sp\else\^{}\fi}\catcode`\%=\active\def%{\%}$\mathdefault{50}$}}%
\end{pgfscope}%
\begin{pgfscope}%
\definecolor{textcolor}{rgb}{0.000000,0.000000,0.000000}%
\pgfsetstrokecolor{textcolor}%
\pgfsetfillcolor{textcolor}%
\pgftext[x=3.127081in,y=0.226595in,,top]{\color{textcolor}{\rmfamily\fontsize{14.000000}{16.800000}\selectfont\catcode`\^=\active\def^{\ifmmode\sp\else\^{}\fi}\catcode`\%=\active\def%{\%}Относительная деформация $\varepsilon$, \%}}%
\end{pgfscope}%
\begin{pgfscope}%
\pgfpathrectangle{\pgfqpoint{0.624162in}{0.530740in}}{\pgfqpoint{5.005838in}{3.099260in}}%
\pgfusepath{clip}%
\pgfsetrectcap%
\pgfsetroundjoin%
\pgfsetlinewidth{0.501875pt}%
\definecolor{currentstroke}{rgb}{0.501961,0.501961,0.501961}%
\pgfsetstrokecolor{currentstroke}%
\pgfsetstrokeopacity{0.500000}%
\pgfsetdash{}{0pt}%
\pgfpathmoveto{\pgfqpoint{0.624162in}{0.530740in}}%
\pgfpathlineto{\pgfqpoint{5.630000in}{0.530740in}}%
\pgfusepath{stroke}%
\end{pgfscope}%
\begin{pgfscope}%
\pgfsetbuttcap%
\pgfsetroundjoin%
\definecolor{currentfill}{rgb}{0.000000,0.000000,0.000000}%
\pgfsetfillcolor{currentfill}%
\pgfsetlinewidth{0.803000pt}%
\definecolor{currentstroke}{rgb}{0.000000,0.000000,0.000000}%
\pgfsetstrokecolor{currentstroke}%
\pgfsetdash{}{0pt}%
\pgfsys@defobject{currentmarker}{\pgfqpoint{-0.048611in}{0.000000in}}{\pgfqpoint{-0.000000in}{0.000000in}}{%
\pgfpathmoveto{\pgfqpoint{-0.000000in}{0.000000in}}%
\pgfpathlineto{\pgfqpoint{-0.048611in}{0.000000in}}%
\pgfusepath{stroke,fill}%
}%
\begin{pgfscope}%
\pgfsys@transformshift{0.624162in}{0.530740in}%
\pgfsys@useobject{currentmarker}{}%
\end{pgfscope}%
\end{pgfscope}%
\begin{pgfscope}%
\definecolor{textcolor}{rgb}{0.000000,0.000000,0.000000}%
\pgfsetstrokecolor{textcolor}%
\pgfsetfillcolor{textcolor}%
\pgftext[x=0.445343in, y=0.472879in, left, base]{\color{textcolor}{\rmfamily\fontsize{12.000000}{14.400000}\selectfont\catcode`\^=\active\def^{\ifmmode\sp\else\^{}\fi}\catcode`\%=\active\def%{\%}$\mathdefault{0}$}}%
\end{pgfscope}%
\begin{pgfscope}%
\pgfpathrectangle{\pgfqpoint{0.624162in}{0.530740in}}{\pgfqpoint{5.005838in}{3.099260in}}%
\pgfusepath{clip}%
\pgfsetrectcap%
\pgfsetroundjoin%
\pgfsetlinewidth{0.501875pt}%
\definecolor{currentstroke}{rgb}{0.501961,0.501961,0.501961}%
\pgfsetstrokecolor{currentstroke}%
\pgfsetstrokeopacity{0.500000}%
\pgfsetdash{}{0pt}%
\pgfpathmoveto{\pgfqpoint{0.624162in}{0.891741in}}%
\pgfpathlineto{\pgfqpoint{5.630000in}{0.891741in}}%
\pgfusepath{stroke}%
\end{pgfscope}%
\begin{pgfscope}%
\pgfsetbuttcap%
\pgfsetroundjoin%
\definecolor{currentfill}{rgb}{0.000000,0.000000,0.000000}%
\pgfsetfillcolor{currentfill}%
\pgfsetlinewidth{0.803000pt}%
\definecolor{currentstroke}{rgb}{0.000000,0.000000,0.000000}%
\pgfsetstrokecolor{currentstroke}%
\pgfsetdash{}{0pt}%
\pgfsys@defobject{currentmarker}{\pgfqpoint{-0.048611in}{0.000000in}}{\pgfqpoint{-0.000000in}{0.000000in}}{%
\pgfpathmoveto{\pgfqpoint{-0.000000in}{0.000000in}}%
\pgfpathlineto{\pgfqpoint{-0.048611in}{0.000000in}}%
\pgfusepath{stroke,fill}%
}%
\begin{pgfscope}%
\pgfsys@transformshift{0.624162in}{0.891741in}%
\pgfsys@useobject{currentmarker}{}%
\end{pgfscope}%
\end{pgfscope}%
\begin{pgfscope}%
\definecolor{textcolor}{rgb}{0.000000,0.000000,0.000000}%
\pgfsetstrokecolor{textcolor}%
\pgfsetfillcolor{textcolor}%
\pgftext[x=0.282151in, y=0.833880in, left, base]{\color{textcolor}{\rmfamily\fontsize{12.000000}{14.400000}\selectfont\catcode`\^=\active\def^{\ifmmode\sp\else\^{}\fi}\catcode`\%=\active\def%{\%}$\mathdefault{100}$}}%
\end{pgfscope}%
\begin{pgfscope}%
\pgfpathrectangle{\pgfqpoint{0.624162in}{0.530740in}}{\pgfqpoint{5.005838in}{3.099260in}}%
\pgfusepath{clip}%
\pgfsetrectcap%
\pgfsetroundjoin%
\pgfsetlinewidth{0.501875pt}%
\definecolor{currentstroke}{rgb}{0.501961,0.501961,0.501961}%
\pgfsetstrokecolor{currentstroke}%
\pgfsetstrokeopacity{0.500000}%
\pgfsetdash{}{0pt}%
\pgfpathmoveto{\pgfqpoint{0.624162in}{1.252743in}}%
\pgfpathlineto{\pgfqpoint{5.630000in}{1.252743in}}%
\pgfusepath{stroke}%
\end{pgfscope}%
\begin{pgfscope}%
\pgfsetbuttcap%
\pgfsetroundjoin%
\definecolor{currentfill}{rgb}{0.000000,0.000000,0.000000}%
\pgfsetfillcolor{currentfill}%
\pgfsetlinewidth{0.803000pt}%
\definecolor{currentstroke}{rgb}{0.000000,0.000000,0.000000}%
\pgfsetstrokecolor{currentstroke}%
\pgfsetdash{}{0pt}%
\pgfsys@defobject{currentmarker}{\pgfqpoint{-0.048611in}{0.000000in}}{\pgfqpoint{-0.000000in}{0.000000in}}{%
\pgfpathmoveto{\pgfqpoint{-0.000000in}{0.000000in}}%
\pgfpathlineto{\pgfqpoint{-0.048611in}{0.000000in}}%
\pgfusepath{stroke,fill}%
}%
\begin{pgfscope}%
\pgfsys@transformshift{0.624162in}{1.252743in}%
\pgfsys@useobject{currentmarker}{}%
\end{pgfscope}%
\end{pgfscope}%
\begin{pgfscope}%
\definecolor{textcolor}{rgb}{0.000000,0.000000,0.000000}%
\pgfsetstrokecolor{textcolor}%
\pgfsetfillcolor{textcolor}%
\pgftext[x=0.282151in, y=1.194882in, left, base]{\color{textcolor}{\rmfamily\fontsize{12.000000}{14.400000}\selectfont\catcode`\^=\active\def^{\ifmmode\sp\else\^{}\fi}\catcode`\%=\active\def%{\%}$\mathdefault{200}$}}%
\end{pgfscope}%
\begin{pgfscope}%
\pgfpathrectangle{\pgfqpoint{0.624162in}{0.530740in}}{\pgfqpoint{5.005838in}{3.099260in}}%
\pgfusepath{clip}%
\pgfsetrectcap%
\pgfsetroundjoin%
\pgfsetlinewidth{0.501875pt}%
\definecolor{currentstroke}{rgb}{0.501961,0.501961,0.501961}%
\pgfsetstrokecolor{currentstroke}%
\pgfsetstrokeopacity{0.500000}%
\pgfsetdash{}{0pt}%
\pgfpathmoveto{\pgfqpoint{0.624162in}{1.613745in}}%
\pgfpathlineto{\pgfqpoint{5.630000in}{1.613745in}}%
\pgfusepath{stroke}%
\end{pgfscope}%
\begin{pgfscope}%
\pgfsetbuttcap%
\pgfsetroundjoin%
\definecolor{currentfill}{rgb}{0.000000,0.000000,0.000000}%
\pgfsetfillcolor{currentfill}%
\pgfsetlinewidth{0.803000pt}%
\definecolor{currentstroke}{rgb}{0.000000,0.000000,0.000000}%
\pgfsetstrokecolor{currentstroke}%
\pgfsetdash{}{0pt}%
\pgfsys@defobject{currentmarker}{\pgfqpoint{-0.048611in}{0.000000in}}{\pgfqpoint{-0.000000in}{0.000000in}}{%
\pgfpathmoveto{\pgfqpoint{-0.000000in}{0.000000in}}%
\pgfpathlineto{\pgfqpoint{-0.048611in}{0.000000in}}%
\pgfusepath{stroke,fill}%
}%
\begin{pgfscope}%
\pgfsys@transformshift{0.624162in}{1.613745in}%
\pgfsys@useobject{currentmarker}{}%
\end{pgfscope}%
\end{pgfscope}%
\begin{pgfscope}%
\definecolor{textcolor}{rgb}{0.000000,0.000000,0.000000}%
\pgfsetstrokecolor{textcolor}%
\pgfsetfillcolor{textcolor}%
\pgftext[x=0.282151in, y=1.555883in, left, base]{\color{textcolor}{\rmfamily\fontsize{12.000000}{14.400000}\selectfont\catcode`\^=\active\def^{\ifmmode\sp\else\^{}\fi}\catcode`\%=\active\def%{\%}$\mathdefault{300}$}}%
\end{pgfscope}%
\begin{pgfscope}%
\pgfpathrectangle{\pgfqpoint{0.624162in}{0.530740in}}{\pgfqpoint{5.005838in}{3.099260in}}%
\pgfusepath{clip}%
\pgfsetrectcap%
\pgfsetroundjoin%
\pgfsetlinewidth{0.501875pt}%
\definecolor{currentstroke}{rgb}{0.501961,0.501961,0.501961}%
\pgfsetstrokecolor{currentstroke}%
\pgfsetstrokeopacity{0.500000}%
\pgfsetdash{}{0pt}%
\pgfpathmoveto{\pgfqpoint{0.624162in}{1.974746in}}%
\pgfpathlineto{\pgfqpoint{5.630000in}{1.974746in}}%
\pgfusepath{stroke}%
\end{pgfscope}%
\begin{pgfscope}%
\pgfsetbuttcap%
\pgfsetroundjoin%
\definecolor{currentfill}{rgb}{0.000000,0.000000,0.000000}%
\pgfsetfillcolor{currentfill}%
\pgfsetlinewidth{0.803000pt}%
\definecolor{currentstroke}{rgb}{0.000000,0.000000,0.000000}%
\pgfsetstrokecolor{currentstroke}%
\pgfsetdash{}{0pt}%
\pgfsys@defobject{currentmarker}{\pgfqpoint{-0.048611in}{0.000000in}}{\pgfqpoint{-0.000000in}{0.000000in}}{%
\pgfpathmoveto{\pgfqpoint{-0.000000in}{0.000000in}}%
\pgfpathlineto{\pgfqpoint{-0.048611in}{0.000000in}}%
\pgfusepath{stroke,fill}%
}%
\begin{pgfscope}%
\pgfsys@transformshift{0.624162in}{1.974746in}%
\pgfsys@useobject{currentmarker}{}%
\end{pgfscope}%
\end{pgfscope}%
\begin{pgfscope}%
\definecolor{textcolor}{rgb}{0.000000,0.000000,0.000000}%
\pgfsetstrokecolor{textcolor}%
\pgfsetfillcolor{textcolor}%
\pgftext[x=0.282151in, y=1.916885in, left, base]{\color{textcolor}{\rmfamily\fontsize{12.000000}{14.400000}\selectfont\catcode`\^=\active\def^{\ifmmode\sp\else\^{}\fi}\catcode`\%=\active\def%{\%}$\mathdefault{400}$}}%
\end{pgfscope}%
\begin{pgfscope}%
\pgfpathrectangle{\pgfqpoint{0.624162in}{0.530740in}}{\pgfqpoint{5.005838in}{3.099260in}}%
\pgfusepath{clip}%
\pgfsetrectcap%
\pgfsetroundjoin%
\pgfsetlinewidth{0.501875pt}%
\definecolor{currentstroke}{rgb}{0.501961,0.501961,0.501961}%
\pgfsetstrokecolor{currentstroke}%
\pgfsetstrokeopacity{0.500000}%
\pgfsetdash{}{0pt}%
\pgfpathmoveto{\pgfqpoint{0.624162in}{2.335748in}}%
\pgfpathlineto{\pgfqpoint{5.630000in}{2.335748in}}%
\pgfusepath{stroke}%
\end{pgfscope}%
\begin{pgfscope}%
\pgfsetbuttcap%
\pgfsetroundjoin%
\definecolor{currentfill}{rgb}{0.000000,0.000000,0.000000}%
\pgfsetfillcolor{currentfill}%
\pgfsetlinewidth{0.803000pt}%
\definecolor{currentstroke}{rgb}{0.000000,0.000000,0.000000}%
\pgfsetstrokecolor{currentstroke}%
\pgfsetdash{}{0pt}%
\pgfsys@defobject{currentmarker}{\pgfqpoint{-0.048611in}{0.000000in}}{\pgfqpoint{-0.000000in}{0.000000in}}{%
\pgfpathmoveto{\pgfqpoint{-0.000000in}{0.000000in}}%
\pgfpathlineto{\pgfqpoint{-0.048611in}{0.000000in}}%
\pgfusepath{stroke,fill}%
}%
\begin{pgfscope}%
\pgfsys@transformshift{0.624162in}{2.335748in}%
\pgfsys@useobject{currentmarker}{}%
\end{pgfscope}%
\end{pgfscope}%
\begin{pgfscope}%
\definecolor{textcolor}{rgb}{0.000000,0.000000,0.000000}%
\pgfsetstrokecolor{textcolor}%
\pgfsetfillcolor{textcolor}%
\pgftext[x=0.282151in, y=2.277887in, left, base]{\color{textcolor}{\rmfamily\fontsize{12.000000}{14.400000}\selectfont\catcode`\^=\active\def^{\ifmmode\sp\else\^{}\fi}\catcode`\%=\active\def%{\%}$\mathdefault{500}$}}%
\end{pgfscope}%
\begin{pgfscope}%
\pgfpathrectangle{\pgfqpoint{0.624162in}{0.530740in}}{\pgfqpoint{5.005838in}{3.099260in}}%
\pgfusepath{clip}%
\pgfsetrectcap%
\pgfsetroundjoin%
\pgfsetlinewidth{0.501875pt}%
\definecolor{currentstroke}{rgb}{0.501961,0.501961,0.501961}%
\pgfsetstrokecolor{currentstroke}%
\pgfsetstrokeopacity{0.500000}%
\pgfsetdash{}{0pt}%
\pgfpathmoveto{\pgfqpoint{0.624162in}{2.696750in}}%
\pgfpathlineto{\pgfqpoint{5.630000in}{2.696750in}}%
\pgfusepath{stroke}%
\end{pgfscope}%
\begin{pgfscope}%
\pgfsetbuttcap%
\pgfsetroundjoin%
\definecolor{currentfill}{rgb}{0.000000,0.000000,0.000000}%
\pgfsetfillcolor{currentfill}%
\pgfsetlinewidth{0.803000pt}%
\definecolor{currentstroke}{rgb}{0.000000,0.000000,0.000000}%
\pgfsetstrokecolor{currentstroke}%
\pgfsetdash{}{0pt}%
\pgfsys@defobject{currentmarker}{\pgfqpoint{-0.048611in}{0.000000in}}{\pgfqpoint{-0.000000in}{0.000000in}}{%
\pgfpathmoveto{\pgfqpoint{-0.000000in}{0.000000in}}%
\pgfpathlineto{\pgfqpoint{-0.048611in}{0.000000in}}%
\pgfusepath{stroke,fill}%
}%
\begin{pgfscope}%
\pgfsys@transformshift{0.624162in}{2.696750in}%
\pgfsys@useobject{currentmarker}{}%
\end{pgfscope}%
\end{pgfscope}%
\begin{pgfscope}%
\definecolor{textcolor}{rgb}{0.000000,0.000000,0.000000}%
\pgfsetstrokecolor{textcolor}%
\pgfsetfillcolor{textcolor}%
\pgftext[x=0.282151in, y=2.638888in, left, base]{\color{textcolor}{\rmfamily\fontsize{12.000000}{14.400000}\selectfont\catcode`\^=\active\def^{\ifmmode\sp\else\^{}\fi}\catcode`\%=\active\def%{\%}$\mathdefault{600}$}}%
\end{pgfscope}%
\begin{pgfscope}%
\pgfpathrectangle{\pgfqpoint{0.624162in}{0.530740in}}{\pgfqpoint{5.005838in}{3.099260in}}%
\pgfusepath{clip}%
\pgfsetrectcap%
\pgfsetroundjoin%
\pgfsetlinewidth{0.501875pt}%
\definecolor{currentstroke}{rgb}{0.501961,0.501961,0.501961}%
\pgfsetstrokecolor{currentstroke}%
\pgfsetstrokeopacity{0.500000}%
\pgfsetdash{}{0pt}%
\pgfpathmoveto{\pgfqpoint{0.624162in}{3.057751in}}%
\pgfpathlineto{\pgfqpoint{5.630000in}{3.057751in}}%
\pgfusepath{stroke}%
\end{pgfscope}%
\begin{pgfscope}%
\pgfsetbuttcap%
\pgfsetroundjoin%
\definecolor{currentfill}{rgb}{0.000000,0.000000,0.000000}%
\pgfsetfillcolor{currentfill}%
\pgfsetlinewidth{0.803000pt}%
\definecolor{currentstroke}{rgb}{0.000000,0.000000,0.000000}%
\pgfsetstrokecolor{currentstroke}%
\pgfsetdash{}{0pt}%
\pgfsys@defobject{currentmarker}{\pgfqpoint{-0.048611in}{0.000000in}}{\pgfqpoint{-0.000000in}{0.000000in}}{%
\pgfpathmoveto{\pgfqpoint{-0.000000in}{0.000000in}}%
\pgfpathlineto{\pgfqpoint{-0.048611in}{0.000000in}}%
\pgfusepath{stroke,fill}%
}%
\begin{pgfscope}%
\pgfsys@transformshift{0.624162in}{3.057751in}%
\pgfsys@useobject{currentmarker}{}%
\end{pgfscope}%
\end{pgfscope}%
\begin{pgfscope}%
\definecolor{textcolor}{rgb}{0.000000,0.000000,0.000000}%
\pgfsetstrokecolor{textcolor}%
\pgfsetfillcolor{textcolor}%
\pgftext[x=0.282151in, y=2.999890in, left, base]{\color{textcolor}{\rmfamily\fontsize{12.000000}{14.400000}\selectfont\catcode`\^=\active\def^{\ifmmode\sp\else\^{}\fi}\catcode`\%=\active\def%{\%}$\mathdefault{700}$}}%
\end{pgfscope}%
\begin{pgfscope}%
\pgfpathrectangle{\pgfqpoint{0.624162in}{0.530740in}}{\pgfqpoint{5.005838in}{3.099260in}}%
\pgfusepath{clip}%
\pgfsetrectcap%
\pgfsetroundjoin%
\pgfsetlinewidth{0.501875pt}%
\definecolor{currentstroke}{rgb}{0.501961,0.501961,0.501961}%
\pgfsetstrokecolor{currentstroke}%
\pgfsetstrokeopacity{0.500000}%
\pgfsetdash{}{0pt}%
\pgfpathmoveto{\pgfqpoint{0.624162in}{3.418753in}}%
\pgfpathlineto{\pgfqpoint{5.630000in}{3.418753in}}%
\pgfusepath{stroke}%
\end{pgfscope}%
\begin{pgfscope}%
\pgfsetbuttcap%
\pgfsetroundjoin%
\definecolor{currentfill}{rgb}{0.000000,0.000000,0.000000}%
\pgfsetfillcolor{currentfill}%
\pgfsetlinewidth{0.803000pt}%
\definecolor{currentstroke}{rgb}{0.000000,0.000000,0.000000}%
\pgfsetstrokecolor{currentstroke}%
\pgfsetdash{}{0pt}%
\pgfsys@defobject{currentmarker}{\pgfqpoint{-0.048611in}{0.000000in}}{\pgfqpoint{-0.000000in}{0.000000in}}{%
\pgfpathmoveto{\pgfqpoint{-0.000000in}{0.000000in}}%
\pgfpathlineto{\pgfqpoint{-0.048611in}{0.000000in}}%
\pgfusepath{stroke,fill}%
}%
\begin{pgfscope}%
\pgfsys@transformshift{0.624162in}{3.418753in}%
\pgfsys@useobject{currentmarker}{}%
\end{pgfscope}%
\end{pgfscope}%
\begin{pgfscope}%
\definecolor{textcolor}{rgb}{0.000000,0.000000,0.000000}%
\pgfsetstrokecolor{textcolor}%
\pgfsetfillcolor{textcolor}%
\pgftext[x=0.282151in, y=3.360891in, left, base]{\color{textcolor}{\rmfamily\fontsize{12.000000}{14.400000}\selectfont\catcode`\^=\active\def^{\ifmmode\sp\else\^{}\fi}\catcode`\%=\active\def%{\%}$\mathdefault{800}$}}%
\end{pgfscope}%
\begin{pgfscope}%
\definecolor{textcolor}{rgb}{0.000000,0.000000,0.000000}%
\pgfsetstrokecolor{textcolor}%
\pgfsetfillcolor{textcolor}%
\pgftext[x=0.226595in,y=2.080370in,,bottom,rotate=90.000000]{\color{textcolor}{\rmfamily\fontsize{14.000000}{16.800000}\selectfont\catcode`\^=\active\def^{\ifmmode\sp\else\^{}\fi}\catcode`\%=\active\def%{\%}Условное напряжение $\sigma$, МПа}}%
\end{pgfscope}%
\begin{pgfscope}%
\pgfpathrectangle{\pgfqpoint{0.624162in}{0.530740in}}{\pgfqpoint{5.005838in}{3.099260in}}%
\pgfusepath{clip}%
\pgfsetrectcap%
\pgfsetroundjoin%
\pgfsetlinewidth{1.505625pt}%
\definecolor{currentstroke}{rgb}{0.839216,0.152941,0.156863}%
\pgfsetstrokecolor{currentstroke}%
\pgfsetdash{}{0pt}%
\pgfpathmoveto{\pgfqpoint{0.624162in}{0.530740in}}%
\pgfpathlineto{\pgfqpoint{0.987944in}{0.553702in}}%
\pgfpathlineto{\pgfqpoint{1.351727in}{0.574388in}}%
\pgfpathlineto{\pgfqpoint{1.715509in}{0.592796in}}%
\pgfpathlineto{\pgfqpoint{2.079292in}{0.608928in}}%
\pgfpathlineto{\pgfqpoint{2.443074in}{0.622783in}}%
\pgfpathlineto{\pgfqpoint{2.806857in}{0.634361in}}%
\pgfpathlineto{\pgfqpoint{3.170639in}{0.643662in}}%
\pgfpathlineto{\pgfqpoint{3.534422in}{0.650686in}}%
\pgfpathlineto{\pgfqpoint{3.898204in}{0.655434in}}%
\pgfpathlineto{\pgfqpoint{4.261987in}{0.657904in}}%
\pgfpathlineto{\pgfqpoint{4.472597in}{0.658294in}}%
\pgfpathlineto{\pgfqpoint{4.510890in}{0.645033in}}%
\pgfpathlineto{\pgfqpoint{4.849099in}{0.520740in}}%
\pgfpathlineto{\pgfqpoint{4.849099in}{0.520740in}}%
\pgfusepath{stroke}%
\end{pgfscope}%
\begin{pgfscope}%
\pgfpathrectangle{\pgfqpoint{0.624162in}{0.530740in}}{\pgfqpoint{5.005838in}{3.099260in}}%
\pgfusepath{clip}%
\pgfsetbuttcap%
\pgfsetroundjoin%
\pgfsetlinewidth{1.505625pt}%
\definecolor{currentstroke}{rgb}{0.549020,0.337255,0.294118}%
\pgfsetstrokecolor{currentstroke}%
\pgfsetdash{{5.550000pt}{2.400000pt}}{0.000000pt}%
\pgfpathmoveto{\pgfqpoint{0.624162in}{0.530740in}}%
\pgfpathlineto{\pgfqpoint{0.786333in}{0.576286in}}%
\pgfpathlineto{\pgfqpoint{0.948504in}{0.619400in}}%
\pgfpathlineto{\pgfqpoint{1.110676in}{0.660080in}}%
\pgfpathlineto{\pgfqpoint{1.272847in}{0.698328in}}%
\pgfpathlineto{\pgfqpoint{1.435018in}{0.734143in}}%
\pgfpathlineto{\pgfqpoint{1.597189in}{0.767526in}}%
\pgfpathlineto{\pgfqpoint{1.759360in}{0.798475in}}%
\pgfpathlineto{\pgfqpoint{1.921532in}{0.826992in}}%
\pgfpathlineto{\pgfqpoint{2.083703in}{0.853075in}}%
\pgfpathlineto{\pgfqpoint{2.245874in}{0.876726in}}%
\pgfpathlineto{\pgfqpoint{2.408045in}{0.897944in}}%
\pgfpathlineto{\pgfqpoint{2.570216in}{0.916730in}}%
\pgfpathlineto{\pgfqpoint{2.732388in}{0.933082in}}%
\pgfpathlineto{\pgfqpoint{2.894559in}{0.947002in}}%
\pgfpathlineto{\pgfqpoint{3.056730in}{0.958488in}}%
\pgfpathlineto{\pgfqpoint{3.218901in}{0.967542in}}%
\pgfpathlineto{\pgfqpoint{3.381072in}{0.974164in}}%
\pgfpathlineto{\pgfqpoint{3.543244in}{0.978352in}}%
\pgfpathlineto{\pgfqpoint{3.705415in}{0.980107in}}%
\pgfpathlineto{\pgfqpoint{3.806772in}{0.979647in}}%
\pgfpathlineto{\pgfqpoint{3.908129in}{0.976787in}}%
\pgfpathlineto{\pgfqpoint{4.009486in}{0.971430in}}%
\pgfpathlineto{\pgfqpoint{4.110843in}{0.963577in}}%
\pgfpathlineto{\pgfqpoint{4.212200in}{0.953227in}}%
\pgfpathlineto{\pgfqpoint{4.313557in}{0.940380in}}%
\pgfpathlineto{\pgfqpoint{4.414914in}{0.925036in}}%
\pgfpathlineto{\pgfqpoint{4.516271in}{0.907196in}}%
\pgfpathlineto{\pgfqpoint{4.617628in}{0.886859in}}%
\pgfpathlineto{\pgfqpoint{4.658171in}{0.878025in}}%
\pgfpathlineto{\pgfqpoint{4.658171in}{0.878025in}}%
\pgfusepath{stroke}%
\end{pgfscope}%
\begin{pgfscope}%
\pgfpathrectangle{\pgfqpoint{0.624162in}{0.530740in}}{\pgfqpoint{5.005838in}{3.099260in}}%
\pgfusepath{clip}%
\pgfsetbuttcap%
\pgfsetroundjoin%
\pgfsetlinewidth{1.505625pt}%
\definecolor{currentstroke}{rgb}{0.090196,0.745098,0.811765}%
\pgfsetstrokecolor{currentstroke}%
\pgfsetdash{{9.600000pt}{2.400000pt}{1.500000pt}{2.400000pt}}{0.000000pt}%
\pgfpathmoveto{\pgfqpoint{0.624162in}{0.530740in}}%
\pgfpathlineto{\pgfqpoint{0.708791in}{3.076447in}}%
\pgfpathlineto{\pgfqpoint{0.722896in}{3.270911in}}%
\pgfpathlineto{\pgfqpoint{0.892156in}{2.846626in}}%
\pgfpathlineto{\pgfqpoint{0.906260in}{2.830035in}}%
\pgfpathlineto{\pgfqpoint{2.091074in}{3.012804in}}%
\pgfpathlineto{\pgfqpoint{2.457802in}{3.122804in}}%
\pgfpathlineto{\pgfqpoint{2.655271in}{3.179783in}}%
\pgfpathlineto{\pgfqpoint{2.824530in}{3.226410in}}%
\pgfpathlineto{\pgfqpoint{2.979685in}{3.266834in}}%
\pgfpathlineto{\pgfqpoint{3.120734in}{3.301291in}}%
\pgfpathlineto{\pgfqpoint{3.247678in}{3.330180in}}%
\pgfpathlineto{\pgfqpoint{3.374622in}{3.356851in}}%
\pgfpathlineto{\pgfqpoint{3.501567in}{3.381115in}}%
\pgfpathlineto{\pgfqpoint{3.614406in}{3.400526in}}%
\pgfpathlineto{\pgfqpoint{3.727246in}{3.417794in}}%
\pgfpathlineto{\pgfqpoint{3.840085in}{3.432823in}}%
\pgfpathlineto{\pgfqpoint{3.952924in}{3.445532in}}%
\pgfpathlineto{\pgfqpoint{4.065764in}{3.455849in}}%
\pgfpathlineto{\pgfqpoint{4.178603in}{3.463718in}}%
\pgfpathlineto{\pgfqpoint{4.291443in}{3.469094in}}%
\pgfpathlineto{\pgfqpoint{4.404282in}{3.471949in}}%
\pgfpathlineto{\pgfqpoint{4.474807in}{3.472446in}}%
\pgfpathlineto{\pgfqpoint{4.488912in}{3.471630in}}%
\pgfpathlineto{\pgfqpoint{4.503016in}{3.469182in}}%
\pgfpathlineto{\pgfqpoint{4.517121in}{3.465102in}}%
\pgfpathlineto{\pgfqpoint{4.531226in}{3.459391in}}%
\pgfpathlineto{\pgfqpoint{4.545331in}{3.452047in}}%
\pgfpathlineto{\pgfqpoint{4.559436in}{3.443072in}}%
\pgfpathlineto{\pgfqpoint{4.573541in}{3.432465in}}%
\pgfpathlineto{\pgfqpoint{4.587646in}{3.420226in}}%
\pgfpathlineto{\pgfqpoint{4.601751in}{3.406355in}}%
\pgfpathlineto{\pgfqpoint{4.629961in}{3.373718in}}%
\pgfpathlineto{\pgfqpoint{4.658171in}{3.334553in}}%
\pgfpathlineto{\pgfqpoint{4.686380in}{3.288861in}}%
\pgfpathlineto{\pgfqpoint{4.714590in}{3.236641in}}%
\pgfpathlineto{\pgfqpoint{4.742800in}{3.177894in}}%
\pgfpathlineto{\pgfqpoint{4.771010in}{3.112620in}}%
\pgfpathlineto{\pgfqpoint{4.799220in}{3.040818in}}%
\pgfpathlineto{\pgfqpoint{4.827430in}{2.962488in}}%
\pgfpathlineto{\pgfqpoint{4.841535in}{2.920876in}}%
\pgfpathlineto{\pgfqpoint{4.841535in}{2.920876in}}%
\pgfusepath{stroke}%
\end{pgfscope}%
\begin{pgfscope}%
\pgfpathrectangle{\pgfqpoint{0.624162in}{0.530740in}}{\pgfqpoint{5.005838in}{3.099260in}}%
\pgfusepath{clip}%
\pgfsetrectcap%
\pgfsetroundjoin%
\pgfsetlinewidth{2.007500pt}%
\definecolor{currentstroke}{rgb}{0.000000,0.000000,0.000000}%
\pgfsetstrokecolor{currentstroke}%
\pgfsetdash{}{0pt}%
\pgfusepath{stroke}%
\end{pgfscope}%
\begin{pgfscope}%
\pgfpathrectangle{\pgfqpoint{0.624162in}{0.530740in}}{\pgfqpoint{5.005838in}{3.099260in}}%
\pgfusepath{clip}%
\pgfsetbuttcap%
\pgfsetmiterjoin%
\definecolor{currentfill}{rgb}{0.000000,0.000000,0.000000}%
\pgfsetfillcolor{currentfill}%
\pgfsetlinewidth{1.003750pt}%
\definecolor{currentstroke}{rgb}{0.000000,0.000000,0.000000}%
\pgfsetstrokecolor{currentstroke}%
\pgfsetdash{}{0pt}%
\pgfsys@defobject{currentmarker}{\pgfqpoint{-0.034722in}{-0.034722in}}{\pgfqpoint{0.034722in}{0.034722in}}{%
\pgfpathmoveto{\pgfqpoint{0.000000in}{0.034722in}}%
\pgfpathlineto{\pgfqpoint{-0.034722in}{-0.034722in}}%
\pgfpathlineto{\pgfqpoint{0.034722in}{-0.034722in}}%
\pgfpathlineto{\pgfqpoint{0.000000in}{0.034722in}}%
\pgfpathclose%
\pgfusepath{stroke,fill}%
}%
\begin{pgfscope}%
\pgfsys@transformshift{5.391627in}{0.321359in}%
\pgfsys@useobject{currentmarker}{}%
\end{pgfscope}%
\end{pgfscope}%
\begin{pgfscope}%
\pgfpathrectangle{\pgfqpoint{0.624162in}{0.530740in}}{\pgfqpoint{5.005838in}{3.099260in}}%
\pgfusepath{clip}%
\pgfsetrectcap%
\pgfsetroundjoin%
\pgfsetlinewidth{2.007500pt}%
\definecolor{currentstroke}{rgb}{0.000000,0.000000,0.000000}%
\pgfsetstrokecolor{currentstroke}%
\pgfsetdash{}{0pt}%
\pgfpathmoveto{\pgfqpoint{4.658171in}{0.878025in}}%
\pgfusepath{stroke}%
\end{pgfscope}%
\begin{pgfscope}%
\pgfpathrectangle{\pgfqpoint{0.624162in}{0.530740in}}{\pgfqpoint{5.005838in}{3.099260in}}%
\pgfusepath{clip}%
\pgfsetbuttcap%
\pgfsetroundjoin%
\definecolor{currentfill}{rgb}{0.000000,0.000000,0.000000}%
\pgfsetfillcolor{currentfill}%
\pgfsetlinewidth{1.003750pt}%
\definecolor{currentstroke}{rgb}{0.000000,0.000000,0.000000}%
\pgfsetstrokecolor{currentstroke}%
\pgfsetdash{}{0pt}%
\pgfsys@defobject{currentmarker}{\pgfqpoint{-0.034722in}{-0.034722in}}{\pgfqpoint{0.034722in}{0.034722in}}{%
\pgfpathmoveto{\pgfqpoint{0.000000in}{-0.034722in}}%
\pgfpathcurveto{\pgfqpoint{0.009208in}{-0.034722in}}{\pgfqpoint{0.018041in}{-0.031064in}}{\pgfqpoint{0.024552in}{-0.024552in}}%
\pgfpathcurveto{\pgfqpoint{0.031064in}{-0.018041in}}{\pgfqpoint{0.034722in}{-0.009208in}}{\pgfqpoint{0.034722in}{0.000000in}}%
\pgfpathcurveto{\pgfqpoint{0.034722in}{0.009208in}}{\pgfqpoint{0.031064in}{0.018041in}}{\pgfqpoint{0.024552in}{0.024552in}}%
\pgfpathcurveto{\pgfqpoint{0.018041in}{0.031064in}}{\pgfqpoint{0.009208in}{0.034722in}}{\pgfqpoint{0.000000in}{0.034722in}}%
\pgfpathcurveto{\pgfqpoint{-0.009208in}{0.034722in}}{\pgfqpoint{-0.018041in}{0.031064in}}{\pgfqpoint{-0.024552in}{0.024552in}}%
\pgfpathcurveto{\pgfqpoint{-0.031064in}{0.018041in}}{\pgfqpoint{-0.034722in}{0.009208in}}{\pgfqpoint{-0.034722in}{0.000000in}}%
\pgfpathcurveto{\pgfqpoint{-0.034722in}{-0.009208in}}{\pgfqpoint{-0.031064in}{-0.018041in}}{\pgfqpoint{-0.024552in}{-0.024552in}}%
\pgfpathcurveto{\pgfqpoint{-0.018041in}{-0.031064in}}{\pgfqpoint{-0.009208in}{-0.034722in}}{\pgfqpoint{0.000000in}{-0.034722in}}%
\pgfpathlineto{\pgfqpoint{0.000000in}{-0.034722in}}%
\pgfpathclose%
\pgfusepath{stroke,fill}%
}%
\begin{pgfscope}%
\pgfsys@transformshift{4.658171in}{0.878025in}%
\pgfsys@useobject{currentmarker}{}%
\end{pgfscope}%
\end{pgfscope}%
\begin{pgfscope}%
\pgfpathrectangle{\pgfqpoint{0.624162in}{0.530740in}}{\pgfqpoint{5.005838in}{3.099260in}}%
\pgfusepath{clip}%
\pgfsetrectcap%
\pgfsetroundjoin%
\pgfsetlinewidth{2.007500pt}%
\definecolor{currentstroke}{rgb}{0.000000,0.000000,0.000000}%
\pgfsetstrokecolor{currentstroke}%
\pgfsetdash{}{0pt}%
\pgfpathmoveto{\pgfqpoint{4.841535in}{2.920876in}}%
\pgfusepath{stroke}%
\end{pgfscope}%
\begin{pgfscope}%
\pgfpathrectangle{\pgfqpoint{0.624162in}{0.530740in}}{\pgfqpoint{5.005838in}{3.099260in}}%
\pgfusepath{clip}%
\pgfsetbuttcap%
\pgfsetmiterjoin%
\definecolor{currentfill}{rgb}{0.000000,0.000000,0.000000}%
\pgfsetfillcolor{currentfill}%
\pgfsetlinewidth{1.003750pt}%
\definecolor{currentstroke}{rgb}{0.000000,0.000000,0.000000}%
\pgfsetstrokecolor{currentstroke}%
\pgfsetdash{}{0pt}%
\pgfsys@defobject{currentmarker}{\pgfqpoint{-0.034722in}{-0.034722in}}{\pgfqpoint{0.034722in}{0.034722in}}{%
\pgfpathmoveto{\pgfqpoint{-0.034722in}{-0.034722in}}%
\pgfpathlineto{\pgfqpoint{0.034722in}{-0.034722in}}%
\pgfpathlineto{\pgfqpoint{0.034722in}{0.034722in}}%
\pgfpathlineto{\pgfqpoint{-0.034722in}{0.034722in}}%
\pgfpathlineto{\pgfqpoint{-0.034722in}{-0.034722in}}%
\pgfpathclose%
\pgfusepath{stroke,fill}%
}%
\begin{pgfscope}%
\pgfsys@transformshift{4.841535in}{2.920876in}%
\pgfsys@useobject{currentmarker}{}%
\end{pgfscope}%
\end{pgfscope}%
\begin{pgfscope}%
\pgfsetrectcap%
\pgfsetmiterjoin%
\pgfsetlinewidth{1.003750pt}%
\definecolor{currentstroke}{rgb}{0.000000,0.000000,0.000000}%
\pgfsetstrokecolor{currentstroke}%
\pgfsetdash{}{0pt}%
\pgfpathmoveto{\pgfqpoint{0.624162in}{0.530740in}}%
\pgfpathlineto{\pgfqpoint{0.624162in}{3.630000in}}%
\pgfusepath{stroke}%
\end{pgfscope}%
\begin{pgfscope}%
\pgfsetrectcap%
\pgfsetmiterjoin%
\pgfsetlinewidth{1.003750pt}%
\definecolor{currentstroke}{rgb}{0.000000,0.000000,0.000000}%
\pgfsetstrokecolor{currentstroke}%
\pgfsetdash{}{0pt}%
\pgfpathmoveto{\pgfqpoint{0.624162in}{0.530740in}}%
\pgfpathlineto{\pgfqpoint{5.630000in}{0.530740in}}%
\pgfusepath{stroke}%
\end{pgfscope}%
\begin{pgfscope}%
\pgfsetbuttcap%
\pgfsetmiterjoin%
\definecolor{currentfill}{rgb}{1.000000,1.000000,1.000000}%
\pgfsetfillcolor{currentfill}%
\pgfsetfillopacity{0.800000}%
\pgfsetlinewidth{1.003750pt}%
\definecolor{currentstroke}{rgb}{0.000000,0.000000,0.000000}%
\pgfsetstrokecolor{currentstroke}%
\pgfsetstrokeopacity{0.800000}%
\pgfsetdash{}{0pt}%
\pgfpathmoveto{\pgfqpoint{4.215352in}{2.789733in}}%
\pgfpathlineto{\pgfqpoint{5.513333in}{2.789733in}}%
\pgfpathquadraticcurveto{\pgfqpoint{5.546667in}{2.789733in}}{\pgfqpoint{5.546667in}{2.823067in}}%
\pgfpathlineto{\pgfqpoint{5.546667in}{3.513333in}}%
\pgfpathquadraticcurveto{\pgfqpoint{5.546667in}{3.546667in}}{\pgfqpoint{5.513333in}{3.546667in}}%
\pgfpathlineto{\pgfqpoint{4.215352in}{3.546667in}}%
\pgfpathquadraticcurveto{\pgfqpoint{4.182018in}{3.546667in}}{\pgfqpoint{4.182018in}{3.513333in}}%
\pgfpathlineto{\pgfqpoint{4.182018in}{2.823067in}}%
\pgfpathquadraticcurveto{\pgfqpoint{4.182018in}{2.789733in}}{\pgfqpoint{4.215352in}{2.789733in}}%
\pgfpathlineto{\pgfqpoint{4.215352in}{2.789733in}}%
\pgfpathclose%
\pgfusepath{stroke,fill}%
\end{pgfscope}%
\begin{pgfscope}%
\pgfsetrectcap%
\pgfsetroundjoin%
\pgfsetlinewidth{1.505625pt}%
\definecolor{currentstroke}{rgb}{0.839216,0.152941,0.156863}%
\pgfsetstrokecolor{currentstroke}%
\pgfsetdash{}{0pt}%
\pgfpathmoveto{\pgfqpoint{4.248685in}{3.421667in}}%
\pgfpathlineto{\pgfqpoint{4.415352in}{3.421667in}}%
\pgfpathlineto{\pgfqpoint{4.582018in}{3.421667in}}%
\pgfusepath{stroke}%
\end{pgfscope}%
\begin{pgfscope}%
\definecolor{textcolor}{rgb}{0.000000,0.000000,0.000000}%
\pgfsetstrokecolor{textcolor}%
\pgfsetfillcolor{textcolor}%
\pgftext[x=4.715352in,y=3.363333in,left,base]{\color{textcolor}{\rmfamily\fontsize{12.000000}{14.400000}\selectfont\catcode`\^=\active\def^{\ifmmode\sp\else\^{}\fi}\catcode`\%=\active\def%{\%}Кожа}}%
\end{pgfscope}%
\begin{pgfscope}%
\pgfsetbuttcap%
\pgfsetroundjoin%
\pgfsetlinewidth{1.505625pt}%
\definecolor{currentstroke}{rgb}{0.549020,0.337255,0.294118}%
\pgfsetstrokecolor{currentstroke}%
\pgfsetdash{{5.550000pt}{2.400000pt}}{0.000000pt}%
\pgfpathmoveto{\pgfqpoint{4.248685in}{3.186022in}}%
\pgfpathlineto{\pgfqpoint{4.415352in}{3.186022in}}%
\pgfpathlineto{\pgfqpoint{4.582018in}{3.186022in}}%
\pgfusepath{stroke}%
\end{pgfscope}%
\begin{pgfscope}%
\definecolor{textcolor}{rgb}{0.000000,0.000000,0.000000}%
\pgfsetstrokecolor{textcolor}%
\pgfsetfillcolor{textcolor}%
\pgftext[x=4.715352in,y=3.127689in,left,base]{\color{textcolor}{\rmfamily\fontsize{12.000000}{14.400000}\selectfont\catcode`\^=\active\def^{\ifmmode\sp\else\^{}\fi}\catcode`\%=\active\def%{\%}Веревка}}%
\end{pgfscope}%
\begin{pgfscope}%
\pgfsetbuttcap%
\pgfsetroundjoin%
\pgfsetlinewidth{1.505625pt}%
\definecolor{currentstroke}{rgb}{0.090196,0.745098,0.811765}%
\pgfsetstrokecolor{currentstroke}%
\pgfsetdash{{9.600000pt}{2.400000pt}{1.500000pt}{2.400000pt}}{0.000000pt}%
\pgfpathmoveto{\pgfqpoint{4.248685in}{2.950378in}}%
\pgfpathlineto{\pgfqpoint{4.415352in}{2.950378in}}%
\pgfpathlineto{\pgfqpoint{4.582018in}{2.950378in}}%
\pgfusepath{stroke}%
\end{pgfscope}%
\begin{pgfscope}%
\definecolor{textcolor}{rgb}{0.000000,0.000000,0.000000}%
\pgfsetstrokecolor{textcolor}%
\pgfsetfillcolor{textcolor}%
\pgftext[x=4.715352in,y=2.892044in,left,base]{\color{textcolor}{\rmfamily\fontsize{12.000000}{14.400000}\selectfont\catcode`\^=\active\def^{\ifmmode\sp\else\^{}\fi}\catcode`\%=\active\def%{\%}Проволока}}%
\end{pgfscope}%
\end{pgfpicture}%
\makeatother%
\endgroup%

    \caption{Диаграммы условных напряжений для исследованных образцов}
    \label{fig:stress_strain}
\end{figure}

\section*{ \boxed{\text{ Задание 6. }} }
\begin{quote}
    \textit{Формулировка вывода о типе разрушения (вязкое/хрупкое) и соответствии материала классу сталей.}
\end{quote}

На основе анализа диаграмм растяжения и расчетных характеристик пластичности ($\delta$ и $\psi$) можно сделать следующие выводы о поведении материалов при одноосном растяжении.

\begin{itemize}
    \item \textit{Кожа.} Образец демонстрирует значительное относительное удлинение ($\delta \approx \luaexec{format_float(results.leather.delta)}\%$) и сужение ($\psi \approx \luaexec{format_float(results.leather.psi)}\%$). Разрушение происходит после существенной пластической деформации. Такое поведение характеризуется как вязкое (пластичное).
    
    \item \textit{Веревка.} Материал показывает значительное относительное удлинение ($\delta \approx \luaexec{format_float(results.rope.delta)}\%$). После достижения максимальной нагрузки происходит падение несущей способности, но образец продолжает деформироваться до разрыва. Относительное сужение составляет $100\%$, что указывает на то, что материал в месте разрыва распался на волокна. Тип разрушения --- вязкий.
    
    \item \textit{Проволока.} Образец имеет относительное удлинение в ($\delta \approx \luaexec{format_float(results.wire.delta)}\%$). Значительное относительное сужение ($\psi \approx \luaexec{format_float(results.wire.psi)}\%$), указывают на наличие развитой пластической деформации перед разрушением. Это характерно для пластичных металлов, таких как сталь. Разрушение является вязким, с образованием <<шейки>>.
\end{itemize}

\section*{Общий вывод}
Все три исследованных материала продемонстрировали пластичный (вязкий) характер разрушения, который сопровождался значительными остаточными деформациями.